\documentclass[11pt]{article}
\usepackage[margin=1in]{geometry}
\usepackage[T1]{fontenc}
\usepackage[utf8]{inputenc}
\usepackage{lmodern}
\usepackage{enumitem}
\usepackage[hidelinks]{hyperref}

\title{Competitive Programming Bootcamp Day 0\\Speaker Notes}
\author{Ryan Pelchat}
\date{August 2026}

\setlist[itemize]{topsep=2pt,itemsep=2pt,parsep=0pt}
\setlist[enumerate]{topsep=2pt,itemsep=2pt,parsep=0pt}

\begin{document}
\maketitle

\section*{Session Goal}
Day 0 should orient students to the ICPC-style contest environment before the bootcamp turns into algorithm practice. The main point is that competitive programming performance is a combination of technical skill, contest rules, team coordination, and disciplined execution under time pressure.

\section*{Slide 1: Title}
\begin{itemize}
  \item Welcome the class and frame Day 0 as orientation rather than a heavy algorithm lecture.
  \item Emphasize that the bootcamp will teach algorithms, but today answers: What are we training for? How does the contest work? How should a team operate?
\end{itemize}

\section*{Slide 2: Today's Roadmap}
\begin{itemize}
  \item Preview the four requested topics: what ICPC is, the rules, performance skills, and team building.
  \item Tell students that the rules discussion is practical, not legalistic: knowing the rules changes strategy and prevents avoidable mistakes.
\end{itemize}

\section*{Slide 3: What Is ICPC?}
\begin{itemize}
  \item Define ICPC as the International Collegiate Programming Contest.
  \item Explain that students work in teams and submit programs to an automatic judge.
  \item Stress that accepted code is the final product, not a written proof or partial explanation.
  \item Transition: the contest format shapes the way we practice.
\end{itemize}

\section*{Slide 4: What Makes ICPC Different?}
\begin{itemize}
  \item Walk through the diagram: three people, one computer, one clock, automatic judging, and visible contest pressure.
  \item Point out that a strong solo programmer can still struggle if the team cannot communicate or share the keyboard.
  \item Ask: What could go wrong if all three teammates try to solve the same problem silently?
\end{itemize}

\section*{Slide 5: The Contest Loop}
\begin{itemize}
  \item Describe the repeatable loop: read, model, solve, code, test, submit or debug.
  \item Explain that most contest mistakes happen when teams skip a step: coding before understanding constraints, submitting before testing, or debugging without a hypothesis.
  \item Tell students this loop will appear throughout the bootcamp.
\end{itemize}

\section*{Slide 6: Typical ICPC-Style Format}
\begin{itemize}
  \item Cover the common setup: three contestants, one workstation, several problems, around five hours, and judged submissions.
  \item Avoid overpromising exact official details. Mention that regions publish their own allowed languages, materials, and site rules.
  \item If students are preparing for a specific regional contest, assign them to read that regional rule page before the first practice contest.
\end{itemize}

\section*{Slide 7: Scoring and Ranking}
\begin{itemize}
  \item Explain the primary scoring rule: more solved problems beats fewer solved problems.
  \item Then explain penalty time: the accepted time for each solved problem plus penalties for wrong attempts before acceptance.
  \item Give the strategic implication: do not guess randomly, but do not be paralyzed by fear of wrong answers either.
  \item Mention that compile-error treatment and tie-breakers should be checked against local rules.
\end{itemize}

\section*{Slide 8: Scoring Example}
\begin{itemize}
  \item Compute the example slowly: Problem A contributes 18 minutes. Problem B contributes 73 minutes plus 40 minutes for two prior rejected attempts. Problem C contributes nothing because it is unsolved.
  \item Ask students which is better: solving one problem quickly or solving two problems with some penalty. The answer depends on solved count first.
  \item Reinforce that unsolved wrong attempts do not add penalty, but they still consume real time.
\end{itemize}

\section*{Slide 9: Contest Conduct}
\begin{itemize}
  \item Frame rules as preserving fairness: teams should solve with their own minds, approved tools, and official contest channels.
  \item Clarification requests go to judges through the contest system, not to friends, online forums, or instructors.
  \item Make the "when unsure" message explicit: ask a judge or proctor.
  \item For bootcamp practice, define which tools students may use during mock contests.
\end{itemize}

\section*{Slide 10: Skills Needed to Perform}
\begin{itemize}
  \item Explain that performance is broader than algorithm knowledge.
  \item Give a quick example for each row: constraints point toward complexity, testing catches off-by-one errors, strategy prevents wasting an hour on the wrong problem.
  \item Ask students to identify which skill area they currently feel strongest in and which one needs the most work.
\end{itemize}

\section*{Slide 11: Technical Skill Stack}
\begin{itemize}
  \item Break technical preparation into fluency, algorithms, and judgment.
  \item Fluency means students can use their language and standard library without searching for every detail.
  \item Algorithms are the reusable patterns the bootcamp will teach.
  \item Judgment is knowing which pattern applies, why it is correct, and whether it fits the constraints.
\end{itemize}

\section*{Slide 12: Team Skill Stack}
\begin{itemize}
  \item Define driver, navigator, and strategist as temporary roles.
  \item Driver is not "the best programmer"; driver is the person currently responsible for turning the plan into code.
  \item Navigator should actively read code and suggest tests, not passively watch.
  \item Strategist protects the team's time and helps decide whether to continue, switch, or ask for clarification.
\end{itemize}

\section*{Slide 13: First 15 Minutes of a Contest}
\begin{itemize}
  \item Explain that the opening minutes are about information gathering.
  \item Encourage teams to divide the problem set quickly, then regroup and compare difficulty estimates.
  \item Make the first target problem a team decision.
  \item Strong warning: do not let one teammate monopolize the computer before the team has an agreed plan.
\end{itemize}

\section*{Slide 14: Submission Discipline}
\begin{itemize}
  \item Present the checklist as a habit that saves time.
  \item Before coding: say algorithm and complexity out loud.
  \item Before submitting: run samples, tiny cases, and boundary cases.
  \item After rejection: avoid random edits. State a likely cause, test it, then change the code.
\end{itemize}

\section*{Slide 15: Icebreaker - Build the Team Map}
\begin{itemize}
  \item Put students in teams of three. If teams are not fixed yet, use temporary groups.
  \item Time box the activity: 8 minutes discussion, 1 minute share-out.
  \item Circulate and listen for whether students are naming concrete skills rather than vague labels.
  \item Share-out prompt: team name, one strength, one growth area, and their "We work best when..." sentence.
\end{itemize}

\section*{Slide 16: Team Challenge - One Keyboard Protocol}
\begin{itemize}
  \item Use the two-sum task because the algorithm is simple enough that the focus stays on communication.
  \item Enforce the rule that only the driver touches the keyboard.
  \item Watch whether the navigator asks useful questions and whether the strategist notices time.
  \item After 5 minutes, rotate roles even if the team is not finished.
  \item Debrief: What helped? What slowed the team down? Did everyone understand the code?
\end{itemize}

\section*{Slide 17: Team Contract}
\begin{itemize}
  \item Ask teams to write short answers, not essays.
  \item Encourage concrete phrases: "pause and prove it", "two more tests", "switch after 10 minutes with no progress".
  \item Tell teams the contract is a living document. They should update it after practice contests.
\end{itemize}

\section*{Slide 18: What Success Looks Like This Week}
\begin{itemize}
  \item Summarize expected Day 0 outcomes.
  \item Students should leave able to describe the contest, explain scoring, name key rules to verify, and use an initial team protocol.
  \item Connect to Day 1: now that students understand the environment, the next sessions build the problem-solving and algorithmic tools.
\end{itemize}

\section*{Slide 19: Sources Checked}
\begin{itemize}
  \item Mention that ICPC rules change by season and region. These notes use the 2026/27 regional and 2026 world finals pages as references.
  \item Recommended instructor action: before delivering this session for a real regional contest, open the local regional page and update allowed languages, materials, and site restrictions.
\end{itemize}

\section*{Slide 20: Closing}
\begin{itemize}
  \item Close with the central message: learn the rules, build the team, solve carefully.
  \item If time remains, let teams revise their contract after the one-keyboard challenge.
\end{itemize}

\section*{Instructor Checklist}
\begin{itemize}
  \item Print or share the official regional rules for the contest students are targeting.
  \item Prepare a simple two-sum prompt for the one-keyboard challenge.
  \item Decide whether bootcamp mock contests allow internet access, AI tools, editor autocomplete, and local notes.
  \item Collect team names and initial role assignments.
  \item Save team contracts for reference after the first practice contest.
\end{itemize}

\section*{Sources}
\begin{itemize}
  \item ICPC Regional Rules 2026/27: \url{https://icpc.global/regionals/rules}
  \item ICPC World Finals Rules 2026: \url{https://icpc.global/worldfinals/rules/}
  \item ICPC Regionals overview: \url{https://icpc.global/regionals/regionals}
\end{itemize}

\end{document}
